% ===================================================================== Methodology
% Every equation below mirrors src/graphnli_axb/models_tokgraph.py (CoveragePartialLabelClassifier,
% BipartiteRelationalLayer) and src/graphnli_axb/token_graph.py (graph construction).
% If the code changes, update this section.
\section{Methodology}
\label{sec:method}

\subsection{Task formulation}
\label{sec:task}

Given a premise $\Prem$ and a hypothesis $\Hyp$, natural language inference (NLI) asks whether $\Prem$
entails $\Hyp$ (\E), contradicts it (\C), or neither (\N). Some corpora only distinguish entailment from
its absence; we keep a single three-way output space $\mathcal{Y}=\{\E,\N,\C\}$ and evaluate two-way
corpora through the induced binary decision \E{} \emph{vs.} $\neg$\E{} (\cref{sec:protocol}).

Our design starts from a simple observation about the label semantics. A hypothesis is entailed when
\emph{every} part of it is supported by the premise and \emph{no} part of it conflicts with the premise;
it is contradicted as soon as \emph{one} part conflicts; otherwise it is neutral.

This view goes back to
probabilistic lexical entailment~\citep{glickman2005probabilistic}, where $P(\Hyp\mid\Prem)$ is a
product over hypothesis words, and to recent atomic-inference models that predict a label per
hypothesis span and combine them with logical rules~\citep{stacey2022logical,wu2023weakly}.
\model{} implements it over \emph{contextual token states} and grounds the per-word decisions in
\emph{typed lexical knowledge}.

\subsection{The \model{} architecture}
\label{sec:architecture}

\begin{figure*}[t]
  \centering
  \includestandalone[width=\textwidth]{figures/cover_architecture}
  \caption{\model{} architecture. (a)~Content words and logical operators of $\Prem$ and $\Hyp$
  become graph nodes; typed lexical relations $r_{ij}$ link premise and hypothesis nodes.
  (b)~A pretrained encoder reads the pair jointly. (c)~Two bipartite relational attention layers
  refine each node with its evidence on the other side, including a learned \emph{null} target.
  (d)~Every hypothesis node $k$ receives a support $s_k$, a conflict $c_k$ and an importance $\omega_k$.
  (e)~A weakest-link coverage and a noisy-OR conflict define a structured distribution $z$, added to
  the \cls{} logits. \todo{update panel (a): operator scope edges are built but not used by \model.}}
  \label{fig:architecture}
\end{figure*}

\Cref{fig:architecture} gives an overview. \model{} keeps the standard cross-encoder classifier and
adds a structured expert whose output is a log-probability distribution over $\mathcal{Y}$. The expert
has three parts: a token graph built from symbolic resources (\cref{sec:graph}), bipartite relational
attention over contextual states (\cref{sec:bipartite}), and a coverage--conflict head with a logical
aggregation (\cref{sec:head}). The two are fused in logit space (\cref{sec:fusion}).

\subsubsection{Encoder}
\label{sec:encoder}

The pair is encoded jointly as \texttt{[CLS]}~$\Prem$~\texttt{[SEP]}~$\Hyp$~\texttt{[SEP]} by a
pretrained Transformer, yielding contextual states $\vect{h}_t\in\mathbb{R}^{d}$ for every subword $t$
and a pooled state $\vect{h}_{\cls}$ (the pooler output when the encoder has one, otherwise the first
token). The baseline classifier is
\begin{equation}
  \vect{\ell}_{\cls} = \vect{W}_c\,\vect{h}_{\cls} + \vect{b}_c \in \mathbb{R}^{3}.
  \label{eq:cls}
\end{equation}
\model{} is encoder-agnostic; we use RoBERTa~\citep{liu2019roberta} for English,
BERTimbau~\citep{souza2020bertimbau} for Portuguese and XLM-R~\citep{conneau2020unsupervised} for the
cross-lingual setting (\cref{sec:protocol}).

\subsubsection{Token graph}
\label{sec:graph}

The graph is aligned with the encoder's subword sequence, so that symbolic evidence attaches to
contextual states rather than to static word embeddings. A word becomes a node if it is a content word
(not punctuation and not in a language-specific stopword list) or a logical operator; the node is
represented by the state of its \emph{first} subword. We write $\Prem$ and $\Hyp$ also for the sets of
premise and hypothesis nodes.

Each word also receives an operator type $o_i\in\Ops$ from language-specific lexicons for English and
Portuguese, with 13 types: \emph{none}, negation, universal, existential, downward-monotone, restrictive,
conditional, disjunction, conjunction, comparative, modal, temporal and number (e.g.\ \textit{not, no, without} / \textit{não,
nenhum, sem} are negations).

Portuguese \textit{se}/\textit{caso} is tagged as
conditional only at the start of a clause, since it is otherwise a reflexive pronoun or a noun.

Premise and hypothesis are connected by typed lexical relations. For every premise--hypothesis word
pair, a relation $r_{ij}\in\Rel$ is assigned by the first rule that
fires: (i)~\emph{number-equal} or \emph{number-conflict} when both words are numerals; (ii)~\emph{identical}
surface form; (iii)~\emph{antonym}, \emph{hypernym}, \emph{hyponym} or \emph{synonym} from a lexical
resource, looked up over a small set of lemma candidates of each word; (iv)~\emph{same-stem} under a light
suffix stripper; otherwise \emph{none}.

Hypernymy is directional: $r_{ij}=\textit{hypernym}$ means the
hypothesis word $j$ is more general than the premise word $i$ ($i\sqsubset j$), and the reverse edge
$r_{ji}$ is the inverse relation. Stopwords never receive lexical relations.

The resources are WordNet~3.0~\citep{miller1995wordnet} and Wiktionary~\citep{ylonen2022wiktextract}
for English, and PT-LKB~\citep{goncalooliveira2018ptlkb,ptlexicalsemantics}, which merges ten Portuguese
lexical resources, and Wiktionary for Portuguese. Pairs listed with
both hypernym directions in the merged resources are discarded instead of guessed, since a wrong
direction swaps forward and reverse entailment. From WordNet we keep single-word lemmas and the three
most frequent synsets per part of speech, to limit edges from rare senses.

\subsubsection{Bipartite relational attention}
\label{sec:bipartite}

Node states are projected to a smaller space and typed by their operator,
\begin{equation}
  \vect{x}^{(0)}_i = \LN\!\big(\vect{W}_p\,\vect{h}_i + \vect{E}_{\mathrm{op}}[o_i]\big)\in\mathbb{R}^{d_g},
  \label{eq:node}
\end{equation}
with $d_g=256$. Each of $L=2$ layers lets every premise node attend to the hypothesis nodes and every
hypothesis node attend to the premise nodes (cross-sentence pairs only), plus a learned \emph{null}
target $\nullnode$ for words without a counterpart. For a node $i$ and a node $j$ on the other side,
\begin{align}
  e_{ij} &= \frac{(\vect{W}_q\vect{x}_i)^{\top}\vect{W}_k\vect{x}_j}{\sqrt{d_g}} + b[r_{ij}],
  \label{eq:score}\\
  \vect{\alpha}_i &= \operatorname{softmax}\big([\,e_{i1},\dots,e_{iN},\,e_{\nullnode}\,]\big),
  \label{eq:align}\\
  \vect{m}_i &= \sum_j \alpha_{ij}\big(\vect{W}_v\vect{x}_j + \vect{e}[r_{ij}]\big) + \alpha_{i\nullnode}\,\vect{v}_{\nullnode},
  \label{eq:message}\\
  \vect{x}_i &\leftarrow \LN\!\big(\vect{x}_i + \MLP([\vect{x}_i;\vect{m}_i;\vect{x}_i\odot\vect{m}_i;|\vect{x}_i-\vect{m}_i|])\big).
  \label{eq:update}
\end{align}
The relation enters twice: as a scalar bias $b[r]$ on the alignment score, which steers \emph{where} a
word looks, and as a vector $\vect{e}[r]$ added to the message, which tells the receiving word
\emph{what kind} of evidence it found (e.g.\ an antonym rather than a synonym). The bias is initialised
to $2$ for identical, synonym, hypernym, hyponym, antonym and number relations, $1$ for same-stem and
$0$ otherwise, so that lexicon edges are preferred alignments at the start of training; all relation
parameters are learned.

The update in \cref{eq:update} follows the matching features of
ESIM~\citep{chen2017enhanced}; $e_{\nullnode}$ and $\vect{v}_{\nullnode}$ are learned, and layers do not
share weights. Unlike pairwise lexical features~\citep{chen2018neural}, the relations here act on
contextual states and are refined over two rounds of alignment.

\subsubsection{Coverage and conflict head}
\label{sec:head}

After the last layer, each hypothesis node $k\in\Hyp$ is scored independently:
\begin{align}
  s_k &= \sigma(\vect{w}_s^{\top}\vect{x}_k + b_s), &\text{(support)}\\
  c_k &= \sigma(\vect{w}_c^{\top}\vect{x}_k + b_c), &\text{(conflict)}\\
  \omega_k &= \softplus(\vect{w}_\omega^{\top}\vect{x}_k + b_\omega). &\text{(importance)}
\end{align}
$s_k$ is the probability that word $k$ is supported by the premise, $c_k$ the probability that it
conflicts with the premise, and $\omega_k\ge 0$ how much the word matters for the decision (function-like
content words can be down-weighted). Biases are initialised to $b_s=2$ and $b_c=-3$, i.e.\ words start
as supported and non-conflicting.

\subsubsection{Logical aggregation}
\label{sec:aggregation}

The word-level scores are combined by two logical rules. The hypothesis is covered if \emph{all} of its words are supported (weakest link) and is free of
conflict if \emph{none} of its words conflicts (noisy-OR):
\begin{align}
  \log \Pcov &= \textstyle\sum_{k\in\Hyp} \omega_k \log s_k, \label{eq:cov}\\
  \log \Pnc  &= \textstyle\sum_{k\in\Hyp} \omega_k \log (1-c_k). \label{eq:noconf}
\end{align}
With $\omega_k=1$, \cref{eq:cov} is the product of word-level support probabilities of
\citet{glickman2005probabilistic}, and \cref{eq:noconf} is a noisy-OR over word-level conflicts. The
three labels follow the decision tree ``any conflict $\Rightarrow$ \C; otherwise all covered
$\Rightarrow$ \E; otherwise \N'':
\begin{equation}
  \begin{aligned}
    z_{\E} &= \log \Pnc + \log \Pcov,\\
    z_{\N} &= \log \Pnc + \log(1-\Pcov),\\
    z_{\C} &= \log(1-\Pnc).
  \end{aligned}
  \label{eq:z}
\end{equation}
Since $\Pnc\Pcov+\Pnc(1-\Pcov)+(1-\Pnc)=1$, $\vect{z}$ is a normalised log-distribution by
construction; no extra softmax or threshold is needed.

Atomic-inference models apply similar rules with
discrete span labels~\citep{stacey2022logical} or fuzzy logic over aligned phrases~\citep{wu2023weakly};
\model{} differs in that the atoms are contextual token nodes whose alignment is shaped by typed lexical
knowledge (\cref{eq:score,eq:message}), and in that the result is fused with, rather than substituted for,
the sentence-level classifier. The per-word values $(s_k,c_k,\omega_k)$ also serve as an explanation of
each prediction.

\subsubsection{Fusion with the sentence classifier}
\label{sec:fusion}

The structured distribution is added to the \cls{} logits with a learned scale $\lambda$ (initialised to
$1$):
\begin{equation}
  \vect{\ell} = \vect{\ell}_{\cls} + \lambda\,\vect{z}, \qquad \hat{y}=\arg\max_{y}\,\ell_y .
  \label{eq:fusion}
\end{equation}
This is a product of experts: the coverage--conflict expert can overrule the encoder when its evidence is
strong (e.g.\ a confident conflict on one word), and defers to it otherwise, so that phenomena the head
cannot express are still handled by the encoder.

\subsection{Training objective}
\label{sec:training}

We use a set-valued negative log-likelihood
$\Lset(\vect{u};Y) = -\log\sum_{y\in Y}\operatorname{softmax}(\vect{u})_y$ over a set $Y$ of admissible
labels, applied to the fused logits and, as an auxiliary term, to the structured expert alone:
\begin{equation}
  \mathcal{L} = \Lset(\vect{\ell};Y) + \beta\,\Lset(\vect{z};Y), \qquad \beta = 0.15 .
  \label{eq:loss}
\end{equation}
The auxiliary term keeps $\vect{z}$ a meaningful classifier on its own instead of a residual correction.
With a single gold label, $Y=\{y^{*}\}$ and \cref{eq:loss} is the usual cross-entropy, which is the case in
all experiments reported here.

The set formulation also allows training on corpora whose negative class
merges neutral and contradiction (e.g.\ the \textsc{none} label of ASSIN), with $Y=\{\N,\C\}$; we leave
this mixed-corpus training to future work.\todo{keep or drop depending on the final experiments.}

The encoder and the new parameters are optimised with AdamW~\citep{loshchilov2019decoupled} and a
linear warm-up/decay schedule; the graph and head parameters use a learning rate $10\times$ larger than
the encoder's, since they start from scratch.

\subsection{Ablations}
\label{sec:ablations}

To identify which components produce the effect, we compare \model{} with five variants that remove one
component each while keeping everything else, including the auxiliary loss, identical (\cref{fig:ablation-map}):
\begin{description}\itemsep0pt
  \item[w/o lexicon] all relations and operators are set to \emph{none}
    ($r_{ij}\!=\!\textit{none}$, $o_i\!=\!\textit{none}$): the graph and head remain, but alignment is
    learned from the encoder states only.
  \item[w/o graph] $L=0$: the head reads $\vect{x}^{(0)}$ directly; this also removes the relation-typed
    attention, hence the lexicon edges.
  \item[w/o conflict] coverage only: $z_{\E}=\log\Pcov$ and $z_{\N}=z_{\C}=\log(1-\Pcov)+\log\tfrac12$.
  \item[w/o coverage] conflict only: $z_{\C}=\log(1-\Pnc)$ and $z_{\E}=z_{\N}=\log\Pnc+\log\tfrac12$.
  \item[MLP head] the logical aggregation (\cref{sec:aggregation}) is replaced by an MLP over the mean and max
    of the hypothesis node states, followed by a log-softmax; this controls for the capacity of the head.
\end{description}
The \cls{} baseline (\cref{eq:cls} alone) is trained with the same hyperparameters.

\begin{figure*}[t]
  \centering
  \includestandalone[width=\textwidth]{figures/ablation_architecture}
  \caption{What each ablation removes (red dashed boxes), drawn on the architecture of
  \cref{fig:architecture}; the numbers refer to the key below the diagram. The \cls{} baseline keeps only
  the encoder and the top $\ell_{\mathrm{CLS}}$ path.}
  \label{fig:ablation-map}
\end{figure*}

\subsection{Knowledge-graph baselines}
\label{sec:baselines}

We also compare with three knowledge-enhanced NLI architectures, re-implemented on the same encoders
and trained under the same protocol and seeds. K-BERT-style injection~\citep{liu2020kbert} appends the
retrieved relations to the premise as text, with a 160-token window so that they are not truncated.
\todo{add KGAnet and KGNLI references.}

The KGAnet-style model encodes a relation-labelled word graph with a relation-aware graph attention
network, and the KGNLI-style model keeps separate semantic and knowledge branches. All three take their
knowledge from rule-based relations (numbers, operators, scope) and English WordNet, so on Portuguese
text their lexical evidence reduces to word identity.

\subsection{Experimental protocol}
\label{sec:protocol}
% Values below follow configs/experiment_0011.yaml and configs/experiment_0012.yaml.
% Update if the final runs change.

\begin{figure*}[t]
  \centering
  \includestandalone[width=\textwidth]{figures/experimental_design}
  \caption{Experimental design. (a)~Three settings, each with its own encoder, training data and lexicon;
  the cross-lingual setting is run in both directions. (b)~All systems of a setting share the same fixed
  hyperparameters. (c)~Each system is trained with 10 seeds; a run keeps the checkpoint with the best
  development macro-F1 and stops after 3 epochs without improvement (at most 8). (d)~The selected
  checkpoint is evaluated once on every test set of its setting. (e)~Systems are compared with paired
  tests over the seeds.}
  \label{fig:design}
\end{figure*}

\subsubsection{Settings and data}
\label{sec:data}

We evaluate three settings (\cref{fig:design,tab:data}).
\emph{English}: RoBERTa-large~\citep{liu2019roberta} trained on SICK~\citep{marelli2014sick}, tested on
the SICK test set, on the SuperGLUE diagnostic AX-b~\citep{wang2019superglue}, which annotates pairs
with linguistic phenomena (lexical semantics, predicate--argument structure, logic, knowledge), and on
three challenge sets described below.

\emph{Portuguese}: BERTimbau-base~\citep{souza2020bertimbau} trained on InferBR~\citep{inferbr2024}
\todo{check InferBR citation}, tested on the InferBR test set and, out of distribution, on the ASSIN
\citep{fonseca2016assin}, ASSIN~2 \citep{real2020assin2} and SICK-BR~\citep{real2018sickbr} test sets.

\emph{Cross-lingual}: XLM-R-base~\citep{conneau2020unsupervised} trained on SICK or on InferBR and tested on
all nine test sets, with the union of the English and Portuguese lexicons in both directions. Relation
types are shared across languages, so relation parameters learned in one language are applied to edges
of the other.

The English challenge sets each isolate one kind of inference. Breaking NLI~\citep{glockner2018breaking}
changes a single word of an SNLI premise (antonyms, co-hyponyms, hypernyms), a direct test of lexical
knowledge and of the conflict detector.

ConjNLI~\citep{saha2020conjnli} tests conjunctive sentences; we use its labelled development set, since the
test labels are not public. The FraCaS suite~\citep{cooper1996fracas} covers formal semantic
phenomena such as quantifiers, plurals and attitudes; multiple premises are concatenated, and the 11
problems without a definite answer are discarded.

SICK-BR is the Portuguese translation of SICK; its test partition is aligned pair by pair with the SICK
test set, so the same pairs can be compared in both languages. ASSIN~2 was built from SICK-BR, and about
36\% of its test pairs (870 of 2{,}448) are translations of SICK training or trial pairs. Models trained on
SICK are therefore scored on the remaining 1{,}578 pairs of ASSIN~2.

\begin{table}[t]
  \centering\small
  \setlength{\tabcolsep}{4pt}
  \begin{tabular}{@{}l l r r r l@{}}
    \toprule
    Dataset & Lang. & Train & Dev & Test & Labels \\
    \midrule
    SICK         & EN & 4{,}439 & 495 & 4{,}906 & 3-way \\
    AX-b         & EN & -- & -- & 1{,}104 & 2-way \\
    Breaking NLI & EN & -- & -- & 8{,}193 & 3-way \\
    ConjNLI      & EN & -- & -- & 623 & 3-way \\
    FraCaS       & EN & -- & -- & 335 & 3-way \\
    \midrule
    InferBR      & PT & 8{,}190 & 633 & 1{,}705 & 3-way \\
    ASSIN        & PT & -- & -- & 4{,}000 & 2-way$^{\dagger}$ \\
    ASSIN~2      & PT & -- & -- & 2{,}448$^{\ddagger}$ & 2-way \\
    SICK-BR      & PT & -- & -- & 4{,}906 & 3-way \\
    \bottomrule
  \end{tabular}
  \caption{Datasets. Only SICK and InferBR are used for training; all other sets are used for testing
  only. $^{\dagger}$Paraphrase pairs are counted as entailment. $^{\ddagger}$1{,}578 pairs for models
  trained on SICK (pairs translated from SICK training or trial data removed).}
  \label{tab:data}
\end{table}

\subsubsection{Training details}
\label{sec:hyper}

Hyperparameters are fixed in advance and shared by all variants of a setting; no per-variant tuning is
performed, so that differences reflect architecture rather than search budget. We train for at most 8
epochs with batch size 32, weight decay $3.5\times10^{-4}$, dropout $0.103$, gradient clipping at $1.0$,
and select the checkpoint with the best development macro-F1, with early stopping after 3 epochs without
improvement.

The learning rate is $10^{-5}$ with 6\% warm-up for RoBERTa-large, which is unstable at larger rates
on small datasets \citep{mosbach2021stability}, $2.04\times10^{-5}$ with 2.75\% warm-up for BERTimbau,
and $2\times10^{-5}$ with 6\% warm-up for XLM-R. \todo{max lengths;
gradient accumulation for RoBERTa-large (2$\times$16, same effective batch).}

\subsubsection{Evaluation}
\label{sec:evaluation}

Our main metric, for every test set, is the Matthews correlation coefficient (MCC) of the binary
decision \E{} \emph{vs.} $\neg$\E{}, obtained by merging \N{} and \C{} predictions. It is the only metric
defined for all our test sets: AX-b, ASSIN and ASSIN~2 only distinguish entailment from non-entailment,
and MCC is the official AX-b metric. It is also robust to the label imbalance of these sets.

The binary view is also required for Breaking NLI: models trained on SICK label most of its
contradictions as neutral, following the annotation convention of SICK, so only the entailment decision
is comparable across training sets. For the in-distribution test sets (SICK and InferBR) we additionally
report three-way macro-F1 in the text.

We further break AX-b results down by the coarse categories of the diagnostic and their fine-grained
phenomena. With 18--140 pairs per phenomenon, these breakdowns are exploratory.

Every system is trained with the same 10 seeds (1001--1010), and results are paired by seed. The seed
fixes all random number generators at the start of a run; because variants create different numbers of
parameters, initialisation and data order are not identical across variants for the same seed, so pairing
blocks run-level variance only partially.

We compare two systems
with the two-sided exact Wilcoxon signed-rank test on the 10 per-seed differences~\citep{wilcoxon1945individual}
and report the mean paired difference with a 95\% percentile bootstrap interval over seeds. Within the
ablation family, $p$-values are corrected with Holm's method~\citep{holm1979simple}. Rankings of more than
two systems use the Friedman test with the Nemenyi post-hoc test~\citep{demsar2006statistical}.
